\documentclass[letterpaper,12pt]{article}

%opening
\title{ImagePy}
\author{Oscar Stomberg}

\begin{document}

\maketitle
\tableofcontents
\clearpage

\section{Reflections}
\subsection{What It Does}
My program takes a source image, and processes it to create a color palette which accurately reflects the colors found in the image. It then applies the source image as my Linux system's desktop background, as well as writing the color palette to various applications' config files, such as my terminal emulator and my desktop's title bar.

\subsection{Why I Like It}
I like this program because it allows me change my system's color scheme quickly when I change desktop backgrounds, without having to do anything manually. It was also a fun exercise in using python, and I learned a lot during it's development.

\subsection{What Resources I Used}
Some resources I use to create this program are:

\begin{itemize}
	\item PyCharm Community Edition, as my IDE. It has a lot of useful features such as built in VCS and code/formatting checking.
	\item "Python For Kids: A Playful Introduction to Programming" by Jason R. Briggs. This is the original book I used a few years ago to learn the basics of Python.
	\item "Automate the Boring Stuff with Python: Practical Programming for Total Beginners" by Al Sweigart. This is the book I used to learn how to use the image processing library.
	\item I used various websites such as \texttt{https://stackoverflow.com/},\\ \texttt{https://www.w3schools.com/python/},
	\texttt{https://docs.python.org/3/}, and various other obscure websites to learn how to use the various data processing libraries available to do what I wanted to do.
\end{itemize}

\subsection{What Surprised Me}
Something that surprised me during this project was that the hard part turned out not to be creating the grouping algorithm, but getting it to run efficiently enough to run it repeatedly so it could achieve the desired number of colors. It was also more challenging than expected to sort the color palette in a consistent and practical manner.

\subsection{What Else Do I Want To Learn}
This year, I would like to learn a lot more about how to structure my code well and write more efficient code.\textsl{\textsl{}}

\end{document}
